\documentclass[../main.tex]{subfiles}

\begin{document}

\onlyinsubfile{
    \renewcommand{\contentsname}{ОГЛАВЛЕНИЕ}
    \setcounter{tocdepth}{3}
    \tableofcontents
}


\chapter*{\MakeUppercase{Заключение}}
\addcontentsline{toc}{chapter}{Заключение}


В рамках курсового проекта было проведено всестороннее исследование методов представления и генерации молекулярных строк с использованием рекуррентных нейронных сетей и других архитектур глубокого обучения. Работа охватила анализ от базовых концепций, заложивших основу генеративной хемоинформатики, до передовых подходов, определяющих современное состояние области, и перспективных направлений ее будущего развития.

В ходе выполнения работы были получены следующие основные результаты в соответствии с поставленными задачами:

\begin{enumerate}
    \item Проанализированы основы представления молекул для задач машинного обучения. Установлено, что нотация SMILES, благодаря своей текстовой природе, стала стандартом де-факто, позволившим применить мощные модели из области обработки естественного языка. Вместе с тем выявлены ее ключевые недостатки: проблема неканоничности, генерация синтаксически некорректных строк и семантическая неадекватность для представления сложных химических концепций.

    \item Рассмотрены базовые генеративные модели. Показано, что рекуррентные нейронные сети (RNN) стали первыми архитектурами, доказавшими принципиальную возможность генерации новых, валидных и химически осмысленных молекул. В свою очередь, вариационные автоэнкодеры (VAE) ввели концепцию непрерывного латентного пространства, что открыло путь к контролируемой генерации и оптимизации свойств молекул путем навигации в этом пространстве.

    \item Проведен детальный сравнительный анализ двух основных парадигм целевой генерации. Подходы на основе VAE и гетероэнкодеров реализуют непрямой контроль через структурирование латентного пространства, что обеспечивает высокое разнообразие генерируемых структур. Методы обучения с подкреплением (RL) осуществляют прямой контроль, напрямую максимизируя целевую функцию (вознаграждение), что обеспечивает исключительную гибкость, но сопряжено с риском «коллапса мод» и требует более сложных техник для поддержания разнообразия.

    \item Исследованы проблемы оценки и ограничения существующих подходов. Подчеркнута важность стандартизированных платформ, таких как MOSES, и комплексных метрик (например, FCD) для объективного сравнения моделей. Определены ключевые тенденции, которые будут формировать будущее генеративного дизайна молекул:
    \begin{itemize}
        \item \textbf{Переход к более мощным архитектурам:} от RNN к Трансформерам и большим языковым моделям (LLM).
        \item \textbf{Переход к более естественным представлениям:} от одномерных строк (SMILES) к графовым нейронным сетям (GNN), работающим напрямую с топологией молекулы.
        \item \textbf{Переход от 2D к 3D:} разработка моделей, способных генерировать не только граф связности, но и пространственную конформацию молекулы, что критически важно для предсказания биологической активности.
    \end{itemize}
\end{enumerate}

Таким образом, все поставленные задачи были выполнены, что позволило в полной мере достичь цели курсового проекта -- провести комплексное исследование и анализ методов нейросетевой генерации молекул.

Результаты работы имеют практическую значимость, поскольку представленный анализ и систематизация подходов могут быть использованы специалистами в области хемоинформатики и дизайна лекарств для выбора наиболее подходящих инструментов и архитектур при решении задач \textit{de novo} генерации молекул с заданными свойствами.

Дальнейшие исследования в этой области могут быть направлены на практическую реализацию и сравнение моделей, основанных на передовых подходах, таких как диффузионные модели для генерации 3D-структур или гибридные системы, объединяющие преимущества обучения с подкреплением и структурированных латентных пространств для достижения одновременно гибкой и разнообразной целевой генерации.

\end{document}